\uuid{0JjE}
\exo7id{3686}
\titre{exo7 3686}
\auteur{quercia}
\organisation{exo7}
\datecreate{2010-03-11}
\isIndication{false}
\isCorrection{false}
\chapitre{Espace euclidien, espace normé}
\sousChapitre{Produit scalaire, norme}
\cpge{mpsi-mp2i.espaces-prehilbertiens.a, pcsi.espaces-prehilbertiens.a, ptsi.geometrie-plan.b, mpsi-mp2i.espaces-prehilbertiens.c, pcsi.espaces-prehilbertiens.c, mpsi-mp2i.changements-bases.a, pcsi.matrices-applications-lineaires.c, ptsi.matrices-applications-lineaires.c, bcpst1.geometrie.f, psi.espaces-prehilbertiens.a, pt.structure-prehilbertienne.a, psi.espaces-prehilbertiens.b, pt.structure-prehilbertienne.c, bcpst2.produit-scalaire.a}
\module{Algèbre}
\niveau{L2}
\difficulte{}

\contenu{
\texte{
Soit $E$ un espace vectoriel euclidien et $\vec u_1,\dots,\vec u_n$ une famille de vecteurs vérifiant :
$\forall\ i\ne j,\ (\vec u_i \mid  \vec u_j) < 0$.
}
\begin{enumerate}
    \item \question{On suppose $(\vec u_1,\dots,\vec u_n)$ libre. Soit $(\vec e_1,\dots,\vec e_n)$ la famille de Schmidt associée et $M$ la matrice de passage de $(\vec u_1,\dots,\vec u_n)$ à $(\vec e_1,\dots,\vec e_n)$. Montrer que $M$ est à coefficients positifs.}
    \item \question{Dans le cas général, démontrer par récurrence sur $n$ que $\mathrm{rg}(\vec u_1,\dots,\vec u_n) \ge n-1$.}
    \item \question{Si $\mathrm{rg}(\vec u_1,\dots,\vec u_n) = n-1$, démontrer que toute famille de $n-1$ vecteurs extraite
    de $(\vec u_1,\dots,\vec u_n)$ est libre, et que les composantes dans cette famille
    du vecteur retiré sont strictement négatives.}
\end{enumerate}
}
